% Part: first-order-logic
% Chapter: syntax-and-semantics
% Section: substitution

\documentclass[../../../include/open-logic-section]{subfiles}

\begin{document}

\olfileid[fr]{fol}{syn}{sub}

\olsection{Substitution}

\begin{defn}[Substitution dans un terme]
Nous définissons récursivement $\Subst{s}{t}{x}$, le résultat de la
\emph{substitution} de~$t$ à chaque occurrence de~$x$ dans~$s$ :
\begin{enumerate}
\item \indcase{s}{c}{$\Subst{\indfrm}{t}{x}$ vaut simplement~$s$.}

\item \indcase{s}{y}{$\Subst{\indfrm}{t}{x}$ vaut également~$s$,
  à condition que $y$ soit une variable et que $y \not\ident x$.}

\item \indcase{s}{x}{$\Subst{\indfrm}{t}{x}$ vaut~$t$.}

\item\indcase{s}{\Atom{f}{t_1, \dots, t_n}}{$\Subst{\indfrmp}{t}{x}$
  vaut $\Atom{f}{\Subst{t_1}{t}{x}, \dots, \Subst{t_n}{t}{x}}$.}
\end{enumerate}
\end{defn}

\begin{defn}
Un terme~$t$ est \emph{!!{free for}} $x$ dans~$!A$ si aucune occurrence
libre de~$x$ dans~$!A$ ne se trouve dans la portée d'un quantificateur
qui lie une variable présente dans~$t$.
\end{defn}

\begin{ex} ~
\begin{enumerate}
\item $\Obj v_8$ est libre pour $\Obj v_1$ dans
  $\lexists[\Obj v_3]\Atom{\Obj A^2_4}{\Obj v_3,\Obj v_1}$.

\item $\Obj f^2_1(\Obj v_1, \Obj v_2)$ n'est \emph{pas} libre pour
  $\Obj v_0$ dans
  $\lforall[\Obj v_2]\Atom{\Obj A^2_4}{\Obj v_0,\Obj v_2}$.
\end{enumerate}
\end{ex}

\begin{defn}[Substitution dans !!a{formula}]
Si $!A$ est !!a{formula}, $x$ !!a{variable} et $t$ un terme
!!{free for}~$x$ dans~$!A$, alors $\Subst{!A}{t}{x}$ est le résultat
de la substitution de~$t$ à toutes les occurrences libres de~$x$
dans~$!A$.
\begin{enumerate}
\tagitem{prvFalse}{\indcase{!A}{\lfalse}{$\Subst{\indfrm}{t}{x}$
    vaut $\lfalse$.}}{}

\tagitem{prvTrue}{\indcase{!A}{\ltrue}{$\Subst{\indfrm}{t}{x}$
    vaut $\ltrue$.}}{}

\item \indcase{!A}{\Atom{P}{t_1,\dots,
    t_n}}{$\Subst{\indfrm}{t}{x}$ vaut
  $\Atom{P}{\Subst{t_1}{t}{x}, \dots, \Subst{t_n}{t}{x}}$.}

\item \indcase{!A}{\eq[t_1][t_2]}{$\Subst{\indfrmp}{t}{x}$ vaut
  $\Subst{t_1}{t}{x} = \Subst{t_2}{t}{x}$.}

\tagitem{prvNot}{\indcase{!A}{\lnot !B}{$\Subst{\indfrmp}{t}{x}$
    vaut $\lnot \Subst{!B}{t}{x}$.}}{}

\tagitem{prvAnd}{\indcase{!A}{(!B \land
    !C)}{$\Subst{\indfrmp}{t}{x}$ vaut $(\Subst{!B}{t}{x} \land
    \Subst{!C}{t}{x})$.}}{}

\tagitem{prvOr}{\indcase{!A}{(!B \lor
    !C)}{$\Subst{\indfrmp}{t}{x}$ vaut $(\Subst{!B}{t}{x} \lor
    \Subst{!C}{t}{x})$.}}{}

\tagitem{prvIf}{\indcase{!A}{(!B \lif
    !C)}{$\Subst{\indfrmp}{t}{x}$ vaut $(\Subst{!B}{t}{x} \lif
    \Subst{!C}{t}{x})$.}}{}

\tagitem{prvIff}{\indcase{!A}{(!B \liff
    !C)}{$\Subst{\indfrmp}{t}{x}$ vaut $(\Subst{!B}{t}{x} \liff
    \Subst{!C}{t}{x})$.}}{}

\tagitem{prvAll}{
\indcase{!A}{\lforall[y][!B]}{$\Subst{\indfrmp}{t}{x}$
    vaut $\lforall[y][\Subst{!B}{t}{x}]$ si $y$ est une variable
    distincte de~$x$ ; sinon, $\Subst{\indfrmp}{t}{x}$
    vaut simplement $\indfrm$.}}{}

\tagitem{prvEx}{
\indcase{!A}{\lexists[y][!B]}{$\Subst{\indfrmp}{t}{x}$
    vaut $\lexists[y][\Subst{!B}{t}{x}]$ si $y$ est une variable
    distincte de~$x$ ; sinon, $\Subst{\indfrmp}{t}{x}$
    vaut simplement $\indfrm$.}}{}
\end{enumerate}
\end{defn}

\begin{explain}
La substitution peut être sans effet : si $x$ n'apparaît pas du tout
dans~$!A$, alors $\Subst{!A}{t}{x}$ vaut simplement~$!A$.

La condition selon laquelle $t$ doit être !!{free for}~$x$ dans~$!A$
exclut notamment le cas suivant. Si
$!A \ident \lexists[y][x < y]$ et $t \ident y$, alors
$\Subst{!A}{t}{x}$ donnerait $\lexists[y][y < y]$. La variable~$y$,
initialement libre, serait « capturée » par le quantificateur
$\lexists[y]$, ce qui est indésirable. Nous aimerions, par exemple,
que si $\lforall[x][!B]$ est vraie, il en soit de même de
$\Subst{!B}{t}{x}$. Or, considérons
$\lforall[x][\lexists[y][x < y]]$, où $!B$ est
$\lexists[y][x < y]$. Cet énoncé est vrai, notamment dans les nombres
naturels : pour tout nombre~$x$, il existe un nombre~$y$ plus grand
que lui. Si nous admettions la substitution de~$y$ à~$x$, nous
obtiendrions $\Subst{!B}{y}{x} \ident \lexists[y][y < y]$, qui est
faux. La condition imposée empêche qu'une variable libre de~$t$
devienne liée par un quantificateur de~$!A$.
\end{explain}

Nous employons souvent la convention suivante pour alléger les
notations. Si $!A$ est !!a{formula} dans laquelle la !!{variable}~$x$
peut apparaître libre, nous écrivons aussi~$!A(x)$ pour le signaler.
Lorsque $!A$ et~$x$ sont clairs dans le contexte, et que $t$ est un
terme supposé libre pour~$x$ dans~$!A(x)$, nous abrégeons
$\Subst{!A}{t}{x}$ par~$!A(t)$. Nous pouvons ainsi dire que « $!A(t)$
est une instance de~$\lforall[x][!A(x)]$ ». Cela signifie que si
$!A$ est une !!{formula} quelconque, $x$ une !!{variable} et $t$ un
terme libre pour~$x$ dans~$!A$, alors $\Subst{!A}{t}{x}$ est une
instance de~$\lforall[x][!A]$.

\end{document}
